\documentclass[11pt]{article}
\usepackage[margin=1in]{geometry}
\usepackage{enumitem}
% =============================================================================
% Preambles/header.tex — shared LaTeX/Beamer preamble for this project
%
% Use from a Beamer lecture by prepending:
%     \input{header}
% and compiling with TEXINPUTS=../Preambles:$TEXINPUTS (the /compile-latex
% skill does this automatically).
%
% PALETTE CONTRACT. Color names defined here MUST match the names in
% Quarto/theme-template.scss so Beamer and Quarto renderings use the same
% palette. The scripts/check-palette-sync.sh script enforces this.
%
% To customize: edit the HEX values below AND the matching $variable in
% Quarto/theme-template.scss, then run scripts/check-palette-sync.sh.
% =============================================================================

% -----------------------------------------------------------------------------
% Engine detection + encoding
%
% Our /compile-latex skill uses XeLaTeX, where inputenc is unnecessary and
% can produce encoding warnings. Only load inputenc under pdfTeX.
% -----------------------------------------------------------------------------
\usepackage{iftex}
\ifPDFTeX
  \usepackage[utf8]{inputenc}
\fi

\usepackage{amsmath, amssymb, amsthm}
\usepackage{booktabs}
\usepackage{graphicx}

% -----------------------------------------------------------------------------
% PALETTE — mirrors Quarto/theme-template.scss variable names.
% Load xcolor and define colors BEFORE anything that references them
% (notably hyperref, hypersetup, and the Beamer color assignments).
% -----------------------------------------------------------------------------
\usepackage{xcolor}

\definecolor{primary-blue}{HTML}{012169}      % headings, accents
\definecolor{primary-gold}{HTML}{B9975B}      % emphasis, borders
\definecolor{highlight-yellow}{HTML}{F2A900}  % markers, alerts
\definecolor{light-bg}{HTML}{E8EDF5}          % title / section backgrounds
\definecolor{jet}{HTML}{1A1A1A}               % body text

% Semantic colors (used in diagrams and annotations)
\definecolor{positive}{HTML}{15803D}          % good / observed / identified
\definecolor{negative}{HTML}{B91C1C}          % bad / problematic / bias
\definecolor{neutral}{HTML}{525252}           % reference / context

% Highlight accents (match .hi-* classes in the SCSS)
\definecolor{hi-slate}{HTML}{314F4F}
\definecolor{hi-green}{HTML}{2E8B57}
\definecolor{hi-red}{HTML}{C41E3A}

% Semantic aliases so skill templates and snippets can write
% \textcolor{accent}{...} without hard-coding "primary-blue".
\colorlet{accent}{primary-blue}
\colorlet{accent2}{primary-gold}

% -----------------------------------------------------------------------------
% Hyperref — loaded AFTER xcolor + palette so link colors resolve.
% -----------------------------------------------------------------------------
\usepackage{hyperref}
\hypersetup{colorlinks=true, linkcolor=primary-blue, urlcolor=primary-blue,
            citecolor=primary-blue, pdfborder={0 0 0}}

% -----------------------------------------------------------------------------
% Course code — set once per deck (after \input{header}, before \begin{document})
% via \coursecode{CS401}. Shown in the footer on every slide so a deck's course
% is identifiable without opening the title page. Empty by default.
% -----------------------------------------------------------------------------
\makeatletter
\newcommand{\coursecode}[1]{\gdef\@coursecodetext{#1}}
\gdef\@coursecodetext{}
\makeatother

% -----------------------------------------------------------------------------
% Beamer theme assignments (only apply under Beamer).
%
% We detect Beamer via \@ifclassloaded{beamer}, which is the documented,
% non-internal way. This must be wrapped in \makeatletter/\makeatother
% because \@ifclassloaded contains an @-catcode character.
% -----------------------------------------------------------------------------
\makeatletter
\@ifclassloaded{beamer}{%
  \usefonttheme{professionalfonts}
  \setbeamercolor{structure}{fg=primary-blue}
  \setbeamercolor{frametitle}{fg=primary-blue}
  \setbeamercolor{title}{fg=primary-blue}
  \setbeamercolor{subtitle}{fg=primary-gold}
  \setbeamercolor{author}{fg=primary-blue}
  \setbeamercolor{institute}{fg=neutral}
  \setbeamercolor{date}{fg=primary-gold}
  \setbeamercolor{itemize item}{fg=primary-blue}
  \setbeamercolor{itemize subitem}{fg=primary-gold}
  \setbeamercolor{alerted text}{fg=primary-gold}
  \setbeamercolor{block title}{fg=white, bg=primary-blue}
  \setbeamercolor{block body}{bg=light-bg}
  % Semantic block variants: block=definition/concept (blue), exampleblock=
  % worked example (green/positive), alertblock=key takeaway or caution (gold).
  % Keep to <=2 colored boxes per slide across ALL three variants (INV-7).
  \setbeamercolor{block title example}{fg=white, bg=positive}
  \setbeamercolor{block body example}{bg=light-bg}
  \setbeamercolor{block title alerted}{fg=white, bg=primary-gold}
  \setbeamercolor{block body alerted}{bg=light-bg}

  % Minimal footer: course code (left) + page X / N (right), no navigation clutter
  \setbeamertemplate{navigation symbols}{}
  \setbeamertemplate{footline}{%
    \color{neutral}\footnotesize\hspace{0.7em}\@coursecodetext\hfill
    \insertframenumber/\inserttotalframenumber\hspace{0.7em}\vspace{0.3em}}
}{}
\makeatother

% -----------------------------------------------------------------------------
% TikZ setup — libraries the snippet gallery uses
% -----------------------------------------------------------------------------
\usepackage{tikz}
\usetikzlibrary{arrows.meta, positioning, calc, decorations.pathreplacing,
                fit, shapes.geometric, backgrounds}

% Shared TikZ styles — snippets and hand-written diagrams can pick these up
% via `\begin{tikzpicture}[every node/.style={...}]` or by naming specific
% styles (e.g., `\node[dag-node] (X) at (0,0) {X};`).
\tikzset{
  % Node styles for causal DAGs and similar
  dag-node/.style = {
    draw=primary-blue, fill=light-bg, rounded corners=2pt,
    minimum width=1.2cm, minimum height=0.9cm, align=center, font=\small
  },
  decision-node/.style = {
    draw=primary-gold, fill=white, diamond, aspect=1.8,
    minimum width=1.8cm, minimum height=1.2cm, align=center, font=\small,
    inner sep=1pt
  },
  % Edge styles — observed vs counterfactual
  observed-edge/.style = {-{Latex[length=2.5mm]}, thick, draw=jet},
  counterfactual-edge/.style = {-{Latex[length=2.5mm]}, thick, draw=neutral, dashed},
  confound-edge/.style = {-{Latex[length=2.5mm]}, thick, draw=negative, dashed},
  % Data-point styles for plots
  observed-dot/.style = {fill=primary-blue, draw=primary-blue, circle, inner sep=1.2pt},
  counterfactual-dot/.style = {fill=white, draw=neutral, circle, inner sep=1.2pt}
}

% -----------------------------------------------------------------------------
% Convenience macros
% -----------------------------------------------------------------------------
\newcommand{\muted}[1]{\textcolor{neutral}{#1}}
\newcommand{\key}[1]{\textcolor{primary-gold}{\textbf{#1}}}
\newcommand{\good}[1]{\textcolor{positive}{#1}}
\newcommand{\bad}[1]{\textcolor{negative}{#1}}

% Standout frame macro: full-bleed dark background for section transitions.
% Beamer-only (uses \begin{frame}/\setbeamercolor/\usebeamerfont) -- do not
% call from an article-class file (e.g. Notes/<CODE>/*-notes.tex). \muted,
% \key, \good, \bad, and \coursecode above are plain xcolor/gdef and are
% safe to use from either class; verified when Notes/article-class support
% was added.
% Expand: \transitionslide{Your transition title}
\newcommand{\transitionslide}[1]{%
  \begingroup
  \setbeamercolor{background canvas}{bg=primary-blue}
  \begin{frame}[plain]
    \centering
    \vfill
    {\usebeamerfont{title}\color{white}\Large #1\par}
    \vfill
  \end{frame}
  \endgroup
}

% Section divider frame (Beamer-only), an alternative to \transitionslide with a
% darker two-line style: near-black (jet) background, a small white label line
% above a larger gold title, and a thin gold rule along the bottom edge.
% Expand: \sectiondivider{Section 2}{Pointers}
\newcommand{\sectiondivider}[2]{%
  \begingroup
  \setbeamercolor{background canvas}{bg=jet}
  \begin{frame}[plain]
    \vspace*{3.0cm}
    {\color{white}\ttfamily\large #1\par}
    \vspace{0.5cm}
    {\color{primary-gold}\LARGE #2\par}
    \begin{tikzpicture}[remember picture, overlay]
      \draw[primary-gold, line width=2.5pt]
        ($(current page.south west)+(0,0.1)$) -- ($(current page.south east)+(0,0.1)$);
    \end{tikzpicture}
  \end{frame}
  \endgroup
}

% -----------------------------------------------------------------------------
% Channel watermark — "Concepts That Click" (YouTube).
%
% Article-class files (CompetitiveExam Q/A sets, Notes, Assignments) get a
% light diagonal draft watermark on every page plus a centered footer naming
% the channel. Beamer decks get a subtle TikZ background watermark instead
% (the footline already carries the course code). Centralized here so every
% MCQ PDF is branded without per-file edits.
% -----------------------------------------------------------------------------
\makeatletter
\@ifclassloaded{article}{%
  \usepackage{draftwatermark}
  \SetWatermarkText{Concepts That Click}
  \SetWatermarkScale{1}
  \SetWatermarkFontSize{2cm}
  \SetWatermarkLightness{0.86}
  \SetWatermarkAngle{45}
  \usepackage{fancyhdr}
  \pagestyle{fancy}
  \fancyhf{}
  \fancyfoot[C]{\footnotesize\textcolor{neutral}{Concepts That Click~|~YouTube: \href{https://www.youtube.com/@concepts.thatclick}{@concepts.thatclick}}}
  \renewcommand{\headrulewidth}{0pt}
  \renewcommand{\footrulewidth}{0pt}
  \fancypagestyle{plain}{%
    \fancyhf{}%
    \fancyfoot[C]{\footnotesize\textcolor{neutral}{Concepts That Click~|~YouTube: \href{https://www.youtube.com/@concepts.thatclick}{@concepts.thatclick}}}%
    \renewcommand{\headrulewidth}{0pt}%
    \renewcommand{\footrulewidth}{0pt}%
  }
}{}
\@ifclassloaded{beamer}{%
  \setbeamertemplate{background}{%
    \begin{tikzpicture}[remember picture, overlay]
      \node[rotate=45, opacity=0.055, align=center]
        at (current page.center)
        {\Huge\textcolor{neutral}{Concepts That Click}};
    \end{tikzpicture}%
  }
}{}
\makeatother 
\coursecode{JKSSB}
\renewcommand{\thesection}{40.\arabic{section}}
\newtheorem{example}{Example}[section]
\newenvironment{definitionbox}[1]{%
  \par\noindent\textbf{Definition (#1).}\ \itshape
}{\par\vspace{0.3em}}
\title{Networking Basics --- Detailed Lecture Notes}
\author{JKSSB Computer Awareness}
\date{\today}

\begin{document}
\maketitle

\section{What a network is: nodes, hosts, clients, servers}
A computer network is formed when two or more computers and related devices
are connected so that they can share data and resources. The shared
resources include files, printers, an internet connection, and application
software. On the office desktop, opening a shared folder, printing to a
shared printer, and browsing the web all use the same underlying network.

\begin{definitionbox}{Computer network}
A set of two or more connected computing devices (nodes) that can exchange
data and share resources such as files, printers, and an internet
connection.
\end{definitionbox}

Every participant in a network is called a \textbf{node}. A node can be a
computer, a printer, a switch, or any device with a network address. A
\textbf{host} is a node that can send or receive data over the network, such
as a desktop computer, a laptop, or a network printer. In practice every host
is a node, but the word host is used when the device is an end point that
produces or consumes data.

\begin{center}
\begin{tabular}{p{0.20\textwidth}p{0.70\textwidth}}
\toprule
\textbf{Term} & \textbf{Meaning in this lesson} \\
\midrule
Network & Two or more connected devices that share data and resources. \\
Node & Any participant device in the network. \\
Host & A node (computer, printer) that sends or receives data. \\
Client & The device or program that requests a service. \\
Server & The computer that provides the service and answers requests. \\
\bottomrule
\end{tabular}
\end{center}

The client--server pair is the most examined relationship in this block. The
\textbf{client} is the device (or program) that requests a service --- a web
page, a file, or a print job. The \textbf{server} is the computer that
provides that service: it stores the shared files or pages and answers the
requests of many clients. When the office desktop opens a file stored on the
department's shared machine, the desktop is the client and the shared machine
is the server.

\begin{definitionbox}{Client and server}
The client requests a service; the server stores the shared resource and
provides it in response to client requests.
\end{definitionbox}

\section{Networks by area: PAN, LAN, MAN, WAN}
Networks are classified by the geographical area they cover. Four names must
be learned together with their full forms, because the exam tests both the
expansion and the coverage in the same item.

\begin{center}
\begin{tabular}{p{0.12\textwidth}p{0.30\textwidth}p{0.48\textwidth}}
\toprule
\textbf{Type} & \textbf{Full form} & \textbf{Coverage and example} \\
\midrule
PAN & Personal Area Network & Around one person: phone, laptop, and earphones linked by Bluetooth. \\
LAN & Local Area Network & One building or campus: an office, school, or laboratory network. \\
MAN & Metropolitan Area Network & One city: a bank's branches linked across a town. \\
WAN & Wide Area Network & A country or the world: the Internet is the largest WAN. \\
\bottomrule
\end{tabular}
\end{center}

A \textbf{PAN} (Personal Area Network) covers the smallest area --- the
devices around a single person, such as a phone connected to a laptop and
wireless earphones. A \textbf{LAN} (Local Area Network) covers one building
or campus, such as all the computers of one office floor connected through a
switch. A \textbf{MAN} (Metropolitan Area Network) spans a city, for example
the branches of one organisation linked across a town. A \textbf{WAN} (Wide
Area Network) spans a country or the whole world; the Internet is the largest
example of a WAN.

\begin{definitionbox}{PAN, LAN, MAN, WAN}
PAN covers one person; LAN covers one building or campus; MAN covers one
city; WAN covers a country or the world.
\end{definitionbox}

\begin{example}
A question gives four coverages --- one person, one building, one city, and
the whole world --- and asks for the matching names. Read the area first:
person means PAN, building means LAN, city means MAN, and worldwide means
WAN. The area decides the answer before the full form is even expanded.
\end{example}

\section{WLAN, intranet, Internet, and ISP}
Three further names complete the vocabulary. A \textbf{WLAN} (Wireless Local
Area Network) is a LAN built with wireless signals instead of cables: the
office Wi-Fi through which laptops connect without a cable is a WLAN. An
\textbf{intranet} is a private network inside one organisation, usable only
by its own staff --- for example the company's internal portal. The
\textbf{Internet} is the worldwide public network of connected networks,
open to anyone with access. The two words differ by one prefix and are a
standing trap: intranet is private to one organisation, while the Internet
is public and worldwide.

An \textbf{ISP} (Internet Service Provider) is the company that provides a
customer with access to the Internet. The home or office connects through its
ISP --- over a phone line, cable, fibre, or mobile link --- and the ISP
carries the traffic onward to the Internet.

\begin{definitionbox}{WLAN, intranet, Internet, ISP}
A WLAN is a wireless LAN; an intranet is a private organisational network;
the Internet is the worldwide public network; an ISP (Internet Service
Provider) gives customers access to the Internet.
\end{definitionbox}

\section{Topology: bus, star, ring, mesh, tree}
The word \textbf{topology} means the arrangement by which nodes are connected
in a network. Five topologies are examined: bus, star, ring, mesh, and tree.
Each one decides how data travels and what happens when one cable or device
fails, and the exam repeatedly asks exactly that: which topology has a
central device, and which one fails fully when the main cable breaks.

In a \textbf{bus} topology, all nodes share one single main cable called the
bus. It is simple and cheap for a small network. Data sent by one node
travels along the whole bus and is accepted by the addressed node. Its
weakness follows from its shape: if the main cable breaks, the whole network
fails.

In a \textbf{star} topology, every node connects to one central device (a
hub or a switch), and all data passes through that centre. If one spoke cable
fails, only that node is affected and the rest keep working. But if the
central device itself fails, the whole network stops. The ordinary office LAN
--- one switch with many desktops plugged into it --- is a star, and it is
the most common topology in exam examples.

In a \textbf{ring} topology, each node connects to exactly two neighbours,
forming a closed circle. Data travels around the ring, usually in one
direction, passed from node to node until it reaches its destination. Each
intermediate node repeats the data it is not addressed to. A break at one
point disturbs the whole ring, unless a second (dual) ring provides a backup
path.

In a \textbf{mesh} topology, every node has a direct link to many (in a full
mesh, all) other nodes. It is the most reliable topology: if one path fails,
another path carries the data. It is also the costliest, because it needs the
most cable and the most ports.

In a \textbf{tree} topology, several star networks are joined through a
central backbone, like the branches of a tree. It suits a large campus: each
department is a star, and the stars join at the centre. Its weakness is the
backbone: if the backbone fails, whole branches lose connection.

\begin{center}
\begin{tabular}{p{0.12\textwidth}p{0.30\textwidth}p{0.48\textwidth}}
\toprule
\textbf{Topology} & \textbf{Shape} & \textbf{What one failure does} \\
\midrule
Bus & One shared cable & Main cable breaks: the whole network fails. \\
Star & Centre with spokes & One spoke fails: one node affected; centre fails: all fail. \\
Ring & Closed circle & One break disturbs the ring. \\
Mesh & Many direct links & One path fails: traffic uses another path. \\
Tree & Stars joined on a backbone & Backbone fails: whole branches fail. \\
\bottomrule
\end{tabular}
\end{center}

\begin{definitionbox}{Topology}
The physical or logical arrangement of nodes in a network: bus (one cable),
star (central device), ring (closed circle), mesh (many direct links), and
tree (stars joined on a backbone).
\end{definitionbox}

\begin{example}
An item states, ``In which topology does the failure of the central device
bring down the whole network?'' Trace the table: only the star routes every
message through one centre, so the answer is star. If the stem instead says
the main cable broke and everything stopped, the shape is the single shared
cable, so the answer is bus.
\end{example}

\section{Transmission media: guided and unguided}
\textbf{Transmission media} are the paths that carry data signals from one
device to another. They divide into two families. \textbf{Guided (wired)
media} confine the signal inside a physical cable. \textbf{Unguided
(wireless) media} carry the signal through air or space without a cable.

There are three guided media. \textbf{Twisted-pair cable} consists of pairs
of copper wires twisted together; it is the ordinary telephone and office-LAN
cable --- cheap and suitable for short ranges. \textbf{Coaxial cable} has a
central copper wire with insulation and a metal shield around it; it is used
for cable TV and older LANs and has better shielding than twisted pair.
\textbf{Fibre-optic cable} uses thin glass or plastic strands carrying light
signals; it is the fastest medium with the longest range and is immune to
electrical noise, but it is also the costliest. Ranked by speed and range,
fibre-optic comes first, then coaxial, then twisted pair.

There are four unguided media at this level. \textbf{Radio waves} travel in
all directions and are used for FM radio and for short-range WLAN and
Bluetooth links. \textbf{Microwaves} travel in a straight line and need line
of sight; they are used for point-to-point links between towers.
\textbf{Satellite links} are microwave signals relayed through a satellite;
they cover very long distances but introduce delay. \textbf{Infrared} works
over a very short range, is blocked by walls, and is used for TV remotes.

\begin{definitionbox}{Guided and unguided media}
Guided media carry signals inside a cable (twisted pair, coaxial,
fibre-optic); unguided media carry signals through air or space (radio
waves, microwaves, satellite links, infrared).
\end{definitionbox}

\section{Networking devices and their roles}
Ten devices are examined, and device--role matching is the highest-yield item
type of this lesson. Each device has one role to memorise.

A \textbf{NIC} (Network Interface Card) is the card --- or the built-in
circuit on the motherboard --- that lets a computer physically connect to a
network. Without it the machine has no network port to plug into.

A \textbf{modem} (modulator-demodulator) converts digital computer signals
into analogue signals for a telephone line and back again. It is the device
that connects a home or office to the ISP over a line that cannot carry
digital signals directly.

A \textbf{hub} connects several nodes in a star and broadcasts every
incoming signal to all of its ports; the addressed node accepts the data and
the rest ignore it. A \textbf{switch} also connects nodes in a star, but it
forwards each message only to the addressed port. The switch is therefore
smarter and faster than the hub, because it does not flood every port with
every message.

A \textbf{repeater} regenerates a weak signal so that it can travel further
along a long cable. It works at the signal level only: it does not choose a
path and does not filter traffic. A \textbf{bridge} connects two LAN
segments and filters traffic, passing across only the data meant for the
other side.

A \textbf{router} connects different networks --- for example a home LAN to
the ISP's network --- and chooses the best path for each packet of data. A
\textbf{gateway} connects networks that use different protocols or data
formats, translating between them so the two sides can understand each
other. An \textbf{access point} is the wireless base station that lets Wi-Fi
devices join a wired network. A \textbf{firewall} blocks unauthorised access
while letting legitimate traffic pass; it is a security device that can be
hardware, software, or both.

\begin{center}
\begin{tabular}{p{0.22\textwidth}p{0.68\textwidth}}
\toprule
\textbf{Device} & \textbf{Role (one line)} \\
\midrule
NIC (Network Interface Card) & Connects a computer to the network. \\
Modem & Converts digital to analogue and back; links the site to the ISP. \\
Hub & Broadcasts every incoming signal to all ports. \\
Switch & Forwards each message only to the addressed port. \\
Repeater & Regenerates a weak signal for a longer run. \\
Bridge & Joins two LAN segments and filters traffic between them. \\
Router & Joins different networks and chooses the best path. \\
Gateway & Translates between networks using different protocols. \\
Access point & Lets wireless (Wi-Fi) devices join the wired network. \\
Firewall & Blocks unauthorised access; a security device. \\
\bottomrule
\end{tabular}
\end{center}

\begin{definitionbox}{The four most confused pairs}
Hub sends to all, switch sends to the addressed port; repeater regenerates
the signal, router chooses the path; bridge joins two similar LAN segments,
gateway translates between different protocols; modem connects to the line,
firewall guards the entry.
\end{definitionbox}

\section{Exam retrieval and common confusions}
Six distinctions prevent most errors on this topic.
\begin{enumerate}
  \item The client \textbf{requests} and the server \textbf{provides}. Do not
    swap the roles.
  \item Read the area first: one person is PAN, one building is LAN, one
    city is MAN, worldwide is WAN.
  \item \textbf{Intranet} is private to one organisation; the
    \textbf{Internet} is public and worldwide.
  \item A hub broadcasts to \textbf{all} ports; a switch forwards only to
    the \textbf{addressed} port.
  \item A repeater only \textbf{regenerates} the signal; the router
    \textbf{chooses the path}. A bridge joins two LAN segments; a gateway
    translates between \textbf{different} protocols.
  \item Fibre-optic is the \textbf{fastest} guided medium; twisted pair is
    the slowest of the three cables.
\end{enumerate}

\begin{example}
An item asks which device ``connects different networks and selects the best
path for each packet.'' Walk the device table: the hub and switch stay inside
one network, the repeater only strengthens the signal, and the gateway
translates protocols --- only the router joins different networks and
chooses paths, so the answer is router.
\end{example}

\subsection*{Exercises}
\begin{enumerate}[label=\arabic*.]
  \item Distinguish a node, a host, a client, and a server.
  \item Expand PAN, LAN, MAN, WAN, WLAN, NIC, and ISP, and state the coverage
    of each network type.
  \item Distinguish an intranet from the Internet.
  \item Describe the five topologies and state, for each, what a single
    failure does.
  \item Distinguish guided from unguided transmission media, giving the three
    cables and the four wireless media of this lesson.
  \item Match each of the ten devices of this lesson to its one-line role.
\end{enumerate}

\subsection*{Solutions}
\begingroup\small
\begin{enumerate}[label=\arabic*.]
  \item A node is any participant device in the network. A host is a node
    (computer, printer) that sends or receives data. A client is the device
    or program that requests a service; a server is the computer that
    provides the service and answers client requests.
  \item PAN --- Personal Area Network (around one person); LAN --- Local Area
    Network (one building or campus); MAN --- Metropolitan Area Network (one
    city); WAN --- Wide Area Network (a country or the world); WLAN ---
    Wireless Local Area Network (a LAN built with wireless signals); NIC ---
    Network Interface Card (lets a computer connect to a network); ISP ---
    Internet Service Provider (gives customers Internet access).
  \item An intranet is a private network inside one organisation, usable only
    by its staff. The Internet is the worldwide public network of connected
    networks, open to anyone with access.
  \item Bus: one shared cable; a break in the main cable fails the whole
    network. Star: every node joins one centre; a spoke failure affects one
    node, a centre failure stops all. Ring: closed circle of neighbours; one
    break disturbs the ring. Mesh: many direct links; one failed path leaves
    others. Tree: stars joined on a backbone; a backbone failure drops whole
    branches.
  \item Guided media confine the signal in a cable: twisted-pair (cheap,
    short range), coaxial (shielded, e.g.\ cable TV), and fibre-optic
    (light signals; fastest, longest range, costliest). Unguided media use
    air or space: radio waves (all directions, short range), microwaves
    (straight line, towers), satellite links (long distance, with delay),
    and infrared (very short range, blocked by walls).
  \item NIC connects the computer; modem converts digital--analogue for the
    line to the ISP; hub broadcasts to all ports; switch forwards to the
    addressed port; repeater regenerates a weak signal; bridge joins two LAN
    segments with filtering; router joins different networks and chooses the
    path; gateway translates between different protocols; access point admits
    Wi-Fi devices; firewall blocks unauthorised access.
\end{enumerate}
\endgroup
\end{document}
